\chapter*{TÓM TẮT}
\addcontentsline{toc}{chapter}{TÓM TẮT}

Sự bùng nổ thông tin trực tuyến tiếng Việt đặt ra hai nhu cầu song song: tóm
tắt nhanh cho người đọc lướt và kiểm chứng đáng tin cậy cho người đọc xác minh.
Các hệ tóm tắt dựa trên một mô hình ngôn ngữ lớn (LLM) đơn lẻ vẫn cho kết quả
không đồng đều, trong khi kỹ thuật Mixture-of-Agents (MoA) --- kết hợp nhiều
LLM qua kiến trúc proposer--aggregator hai tầng --- chưa được nghiên cứu cho
tiếng Việt. Đồng thời, các metric đánh giá tóm tắt phổ biến trong nghiên cứu
tiếng Việt (ROUGE, BLEU, BERTScore) đo độ trùng lặp $n$-gram hoặc embedding
với văn bản gốc, không phản ánh chính xác chất lượng tóm tắt khi hệ thống
được phép viết lại.

Khóa luận đóng góp ba kết quả. \textit{Thứ nhất}, xây dựng \textbf{Fiber} ---
một tiện ích trình duyệt và máy chủ Next.js cho phép tóm tắt và kiểm chứng
tin tức tiếng Việt theo thời gian thực, hỗ trợ 14 LLM thuộc OpenAI, Anthropic
và Google. \textit{Thứ hai}, triển khai pipeline Mixture-of-Agents (Wang et
al., 2024) cho bài toán tóm tắt có grounding tiếng Việt, đưa thêm
\textit{kết nối dư} (residual connection) văn bản gốc vào prompt của
aggregator để giữ tính trung thực. \textit{Thứ ba}, đề xuất \textbf{khung
đánh giá ba trục}: trục~A đo độ lưu giữ nội dung (ROUGE, BLEU, BERTScore),
trục~B đo chất lượng tự động qua LLM-judge (rubric, so sánh cặp, kiểm chứng
sự kiện), trục~C đánh giá bởi con người (xếp hạng mù $K$-way và thống kê
đồng thuận). Khung đánh giá này chỉ ra một cách thực nghiệm vì sao các metric
overlap không đủ để đánh giá hệ tóm tắt dạng aggregator.

Trên tập dữ liệu cân bằng 50 bài báo \textit{tienphong.vn} thuộc 5 chủ đề,
pipeline MoA thắng baseline mạnh nhất (GPT-4o đơn lẻ) trong $37/48$ cặp
($77{,}1\%$, sign-test $p = 0{,}0002$); thắng từng proposer với mức ý nghĩa
thống kê ($98{,}0\%$, $83{,}7\%$, $69{,}6\%$ tương ứng với GPT-4o-mini,
Claude Haiku 4.5 và Gemini 2.5 Flash). Cả điểm rubric và metric overlap đều
ủng hộ kết luận này. Nghiên cứu thí điểm với 4 người đánh giá trên toàn bộ
$48$ task (n=192) phát hiện khoảng cách $23$ điểm phần trăm giữa đánh giá
của LLM-judge và con người: người dùng thích bản fusion chỉ $53{,}6\%$ lần,
không phải $77{,}1\%$. Khoảng cách này, cùng tương tác theo chủ đề (người đọc
chuộng bản tóm tắt ngắn cho tin hành chính, chuộng bản fusion cho tin văn hóa
giải trí), là phát hiện phương pháp luận trọng tâm của khóa luận.

\vspace{0.4em}
\textbf{Từ khóa:} tóm tắt tiếng Việt, kiểm chứng tin tức, mô hình ngôn ngữ
lớn, Mixture-of-Agents, LLM-as-judge, tiện ích trình duyệt, đánh giá ba trục.
